% =========================================================
% The Nested Interval Construction of ℝ (Cantor)
% =========================================================

\section{Nested Interval Construction of $\mathbb{R}$ (Cantor)}
\label{sec:reals-nested-intervals}

\begin{tcolorbox}[colback=gray!6, colframe=gray!40, arc=2pt,
  left=6pt, right=6pt, top=4pt, bottom=4pt,
  title={\small\textbf{Cantor Nested Intervals — Quick Reference}},
  fonttitle=\small\bfseries]
\small
\begin{tabular}{l l l}
\toprule
\textbf{Concept} & \textbf{Meaning} & \textbf{Detail} \\
\midrule
Nested intervals & $I_{n+1}\subseteq I_n$ & \hyperref[def:nested-intervals]{↓ Def} \\
Diameter / length & $\operatorname{diam}(I_n)$ shrinks to $0$ & \hyperref[def:diameter]{↓ Def} \\
Cantor intersection principle & $\bigcap_n I_n$ is a single point & \hyperref[thm:cantor-nested]{↓ Thm} \\
Real as ``limit point'' of rationals & coded by a shrinking rational interval & \hyperref[def:real-as-nested]{↓ Def} \\
\bottomrule
\end{tabular}
\end{tcolorbox}

\vspace{1em}

% ---------------------------------------------------------
% Definitions and Theorems
% ---------------------------------------------------------

\begin{tcolorbox}[colback=defbox, colframe=defborder, arc=2pt,
  left=6pt, right=6pt, top=4pt, bottom=4pt,
  title={\small\textbf{Definition (Nested Interval Sequence)}},
  fonttitle=\small\bfseries]
\label{def:nested-intervals}
A \emph{nested interval sequence} is a sequence of nonempty closed intervals
\[
I_n = [a_n,b_n] \subseteq \mathbb{Q}
\]
such that for all $n\in\mathbb{N}$,
\[
I_{n+1}\subseteq I_n
\quad\text{(equivalently, } a_n \le a_{n+1} \le b_{n+1} \le b_n\text{).}
\]
\end{tcolorbox}

\begin{tcolorbox}[colback=defbox, colframe=defborder, arc=2pt,
  left=6pt, right=6pt, top=4pt, bottom=4pt,
  title={\small\textbf{Definition (Diameter / Length)}},
  fonttitle=\small\bfseries]
\label{def:diameter}
For an interval $I=[a,b]$ (with $a\le b$), its \emph{diameter} (length) is
\[
\operatorname{diam}(I) := b-a.
\]
For a nested sequence $I_n=[a_n,b_n]$, we say it \emph{shrinks to a point} if
\[
\operatorname{diam}(I_n) \longrightarrow 0 \quad \text{as } n\to\infty.
\]
\end{tcolorbox}

\begin{tcolorbox}[colback=thmbox, colframe=thmborder, arc=2pt,
  left=6pt, right=6pt, top=4pt, bottom=4pt,
  title={\small\textbf{Theorem (Cantor Nested Interval Principle)}},
  fonttitle=\small\bfseries]
\label{thm:cantor-nested}
Let $(I_n)$ be a nested interval sequence of nonempty closed intervals
\[
I_n=[a_n,b_n] \subseteq \mathbb{R}
\]
such that $\operatorname{diam}(I_n)\to 0$.
Then there exists a unique real number $x\in\mathbb{R}$ such that
\[
x\in I_n \quad\text{for all } n,
\qquad\text{i.e.}\qquad
\bigcap_{n=1}^\infty I_n = \{x\}.
\]
\end{tcolorbox}

\begin{tcolorbox}[colback=defbox, colframe=defborder, arc=2pt,
  left=6pt, right=6pt, top=4pt, bottom=4pt,
  title={\small\textbf{Definition (Real Number via Nested Rational Intervals)}},
  fonttitle=\small\bfseries]
\label{def:real-as-nested}
Let $\mathcal{N}$ be the set of all nested sequences of nonempty closed
\emph{rational} intervals
\[
I_n=[a_n,b_n]\subseteq \mathbb{Q},
\qquad I_{n+1}\subseteq I_n,
\qquad \operatorname{diam}(I_n)\to 0.
\]
Define an equivalence relation $\sim$ on $\mathcal{N}$ by
\[
(I_n)\sim(J_n)
\quad\Longleftrightarrow\quad
\forall n\ \exists m\ \bigl(I_m\subseteq J_n \ \wedge\  J_m\subseteq I_n\bigr).
\]
(Informally: the two sequences eventually refine each other arbitrarily.)
A \emph{real number} is an equivalence class $[(I_n)]$ under $\sim$.

We denote the resulting set of equivalence classes by $\mathbb{R}$.
\end{tcolorbox}

\begin{tcolorbox}[colback=propbox, colframe=propborder, arc=2pt,
  left=6pt, right=6pt, top=4pt, bottom=4pt,
  title={\small\textbf{Remark (Endpoints approximate the same point)}} ,
  fonttitle=\small\bfseries]
If $(I_n)=[a_n,b_n]$ is nested with $\operatorname{diam}(I_n)\to 0$, then
\[
(b_n-a_n)\to 0,
\]
so $(a_n)$ and $(b_n)$ are ``squeezed together.'' In a completed number system,
this forces them to converge to the \emph{same} point $x$, and the intervals
are precisely rational ``traps'' for $x$.
\end{tcolorbox}

% ---------------------------------------------------------
% Consequences
% ---------------------------------------------------------

\subsection{Consequences}

\begin{remark}[How this constructs completeness]
The Cantor principle (\autoref{thm:cantor-nested}) is a completeness statement:
it says that whenever you can trap a number by nested closed intervals whose
lengths go to $0$, there is a point (indeed a unique point) trapped by all of them.

In the rational numbers $\mathbb{Q}$ this fails: there are nested rational
intervals with diameters $\to 0$ whose ``intended'' intersection point is
irrational (e.g.\ intervals shrinking around $\sqrt{2}$), so the intersection
in $\mathbb{Q}$ is empty. The construction above formally adjoins the missing
limit points by \emph{declaring} each shrinking nested rational interval
sequence to represent a real number.
\end{remark}

\begin{remark}[Comparison with Cauchy completion]
This is morally the same as the Cauchy-sequence construction:
\begin{itemize}
  \item A Cauchy sequence is a list of rational approximations that eventually
        get arbitrarily close.
  \item A nested rational interval sequence is a sequence of rational
        \emph{error bars} whose widths go to $0$.
\end{itemize}
One can convert between the two encodings:
choose a point $q_n\in I_n$ to get a Cauchy sequence $(q_n)$, and conversely
use Cauchy tails to build shrinking intervals.
Both yield the same $\mathbb{R}$ up to canonical isomorphism.
\end{remark}

\begin{remark}[Status]
\textit{This section is intentionally focused on the construction idea and the
core nested-interval principle. Full development (addition/multiplication on
classes, order, and the full completeness axiom) can be added after your
Cauchy-section, since the proofs mirror each other.}
\end{remark}
